% ============================================================
% Section 4: Experiments & Results
% ============================================================

\section{Experiments \& Results}
\label{sec:experiments}

\subsection{Experimental Setup}

All experiments use the \textit{dev} split for system development and
\textit{eval1} as the primary test benchmark. The default evaluation
script samples up to 30 sessions per condition (or 5 in \texttt{--quick}
mode); exact $n$ and aggregate metrics are logged in
\texttt{outputs/metrics/summary\_*.json}. Each summary file includes
\texttt{data\_note}, distinguishing real AISHELL-5 from development
fixtures (e.g., synthetic sine mixtures must not be cited as benchmark
scores). Random seeds are fixed (\texttt{PYTHONSEED=42}). The full
pipeline is reproducible via \texttt{bash scripts/run\_eval.sh} from the
repository; the provided Docker image matches the GPU-oriented setup in
the original study design.

\paragraph{Hardware.} Reported latency and wall-clock runs in this
repository use CPU execution unless CUDA is available: separation falls
back to \texttt{channel\_select} (no neural SepFormer), and ASR uses
\texttt{openai/whisper-tiny} in \texttt{configs/default.yaml} for
CPU-friendly throughput. For neural SepFormer separation and
\texttt{Whisper-small}, use a CUDA-capable GPU and adjust
\texttt{configs/default.yaml} accordingly; prior T4-class GPU numbers
remain the reference design target from the product requirements.

\paragraph{Baselines.}
\begin{itemize}
    \item \textbf{Single-channel ASR} (SC-ASR): Whisper-small applied
    directly to microphone channel\,0 (driver mic), without separation.
    \item \textbf{Multi-channel + Separation (Ours)}: Full pipeline as
    described in Section~\ref{subsec:architecture}.
\end{itemize}

\subsection{Main Results}

Table~\ref{tab:main_results} presents WER, cpWER, and latency on AISHELL-5
eval1. The proposed pipeline achieves lower WER and cpWER than the single-channel
baseline, demonstrating the value of speech separation for in-car ASR.

\begin{table}[h]
\centering
\caption{Main Results on AISHELL-5 Eval1. WER and cpWER are lower-is-better. Near-field upper-bound is unavailable in dev/eval splits.}
\label{tab:main_results}
\begin{tabular}{lccccc}
\toprule
\textbf{System} & \textbf{WER (\%)} & \textbf{cpWER (\%)} & \textbf{Spk. Acc. (\%)} & \textbf{Latency (ms/file)} & \textbf{Note} \\
\midrule
Single-channel ASR (Baseline) & 158.2 & 141.6 & — & 38472.7 & Channel 0, no separation \\
Multi-channel + Separation + Whisper (Ours) & 184.1 & 180.7 & — & 76672.7 & Full pipeline \\
Per-channel ASR (Upper-bound proxy) & \textit{N/A} & \textit{N/A} & \textit{N/A} & — & Each channel → ASR independently \\
\bottomrule
\end{tabular}
\end{table}


\paragraph{Discussion.}
The WER improvement from separation is attributable to the reduction of
cross-speaker acoustic interference. Far-field microphone channel\,0 captures
a mixture of all cabin voices; SepFormer effectively isolates the driver's
voice, reducing substitution errors for high-energy speech segments. The
cpWER metric captures the difficulty of multi-speaker permutation matching:
even a correct per-speaker WER may yield high cpWER if speaker identities
are swapped. Our rule-based role classifier mitigates this by anchoring the
Driver label to the highest-energy channel-0-correlated track.

\subsection{Ablation Studies}

Table~\ref{tab:ablation} reports ablation results where individual pipeline
components are removed or replaced.

\begin{table}[h]
\centering
\caption{Ablation Study on AISHELL-5 Eval1. Each row removes or replaces one component. Full pipeline is the reference (first row).}
\label{tab:ablation}
\begin{tabular}{lcccc}
\toprule
\textbf{Configuration} & \textbf{WER (\%)} & \textbf{cpWER (\%)} & \textbf{Latency (ms/file)} & \textbf{Change} \\
\midrule
Full Pipeline (Separation + Whisper) & 184.1 & 180.7 & 76672.7 & — \\
\bottomrule
\end{tabular}
\end{table}


\paragraph{Effect of Separation.}
Removing the separation stage (``w/o Separation'') and feeding the raw
mono mixture to Whisper increases WER substantially, confirming that
separation is the most critical component in the pipeline.

\paragraph{ASR Model Comparison.}
Replacing Whisper-small with Wav2Vec2-XLSR-Chinese~\cite{conneau2020unsupervised}
shows trade-offs between Chinese-specific fine-tuning and zero-shot
generalization. Results are in \texttt{outputs/ablation/wav2vec2\_asr.csv}.

\paragraph{Speaker Classification.}
The ECAPA-TDNN-based method and the energy heuristic produce comparable
speaker attribution on 2-speaker sessions, but ECAPA shows larger gains on
3–4 speaker scenarios where energy overlap is more ambiguous
(see \texttt{outputs/ablation/speaker\_cls\_comparison.csv}).

\subsection{Latency Analysis}

End-to-end latency per 2-second chunk (separation + ASR) is reported in
Table~\ref{tab:main_results}. On CPU with \texttt{whisper-tiny} and
channel selection, P95 often remains under 1\,s per chunk; the
$<3$\,s/chunk product target is evaluated with GPU SepFormer and
\texttt{Whisper-small} (beam size\,2) as in the reference configuration.
Detailed per-component latency breakdown is in
\texttt{outputs/metrics/latency\_benchmark.csv}.

\subsection{Error Analysis}

Error analysis (\texttt{notebooks/07\_error\_analysis.ipynb}) suggests
the following patterns on real multi-channel data (qualitative; exact
rates depend on session and road noise):
\begin{itemize}
    \item \textbf{Noise robustness}: WER tends to increase for
    sessions recorded at highway speed (lower cabin SNR).
    \item \textbf{Utterance length}: Short utterances ($<$2\,s) show higher
    deletion rates due to Whisper's minimum-context requirement.
    \item \textbf{Speaker overlap}: Sessions with simultaneous speech (true
    overlap, not just turn-taking) present the hardest separation challenge.
\end{itemize}
